\documentstyle[%


righttag%


,11pt%


,amssymb%


,russian%               remove in Eng.


,izvru%                 Set izven% in Eng.


% ,izvru-ks             for brief communications


]{amsart}





%


\newcommand{\ind}{\operatorname{ind}}


\newcommand{\col}{\operatorname{col}}


\newcommand{\eqdef}{\overset{\mathrm{def}}{=}{}}


\newcommand {\N}{{\Bbb N}}


\newcommand{\wsn}{WS^n[a,t_1,\dotsc,t_m,b]}


\newcommand{\mes}{\operatorname{mes}}


\newcommand{\shh}{@>{{\cal H}{\cal H}}>>}


\newcommand{\sph}{@>{{\cal P}{\cal H}}>>}


\newcommand{\sh}{@>{\cal H}>>}


\newcommand{\ssp}{@>{\cal P}>>}


\newcommand{\ps}{\cal P}


\newcommand{\xk}{\{x_k\}}


\newcommand{\lkxfk}{\cal L_kx=f_k}


\newcommand{\lkxak}{l_kx=\alpha_k}


\newcommand{\kodi}{k=1,2,\dotsc}


\newcommand{\tts}[1]{}


\renewcommand{\baselinestretch}{0.97}





%\hoffset=-15mm%                  For Eng.


\setcounter{page}{41}


\god{2000}


\nomer{4~(455)} %                        Remove in Eng.


%\nomer{Vol.\,44, No.\,4}%               Set in Eng.


%\str{\pageref{begin}--\pageref{end}}%  Set in Eng.


\udk{517.929}





\author{\it Л.Ф.\,Рахматуллина}


% NB:   ^^ remove in Eng.


\title{Непрерывная зависимость от параметров\\ решения


линейной краевой задачи}


\thanks{Работа выполнена при финансовой поддержке Российского фонда


фундаментальных исследований (96-15-96195, 99-01-01278) и Конкурсного


центра по исследованиям в области


фундаментального естествознания, Санкт-Петербург.}





\date{}


%\pagestyle{myheadings}%         Set in Eng.


\pagestyle{plain} %              Remove in Eng.


\begin{document}


%\thispagestyle{plain}%          Set in Eng.


\maketitle


\label{begin}








В [1] изучались условия $\ps$-сходимости последовательности


$\xk$ решений линейных краевых задач


$$


\lkxfk,\quad \lkxak,\quad \kodi,


$$


к решению $x_0$ ``предельной'' задачи


$$


{\cal L}_0x=f_0,\quad l_0x=\alpha_0.


$$


Здесь ${\cal L}_k:D_k\to B_k$


--- линейные ограниченные н\"eтеровы операторы,


$\ind{\cal L}_k=N$; $l_k:D_k\to R^N$


--- линейные ограниченные вектор-функционалы;


$\alpha_k\in{R}^N$; $D_k$, $B_k$


--- банаховы пространства, причем $D_k$ изоморфно


$B_k\times R^N$, $k=0,1,\dotsc$,


${\cal P} = \{{\cal P}_k\}_{k\in \N}$


--- связывающая система изоморфизмов пространств


$D_0$ и $D_k$.





Ниже общие утверждения работы [1]


реализуются для краевых задач в конкретных пространствах.





Являясь продолжением работы [1],


предлагаемая статья использует обозначения и терминологию этой работы.





\section{Уравнение $n$-го порядка}





Пусть


$\lim\limits_{k\to\infty} a^k =a^0$,


$\lim\limits_{k\to\infty} b^k=b^0$;


$W^n[a^k,b^k]$ --- пространство функций


$x:[a^k,b^k]\to R^1$, имеющих абсолютно непрерывные производные до


$(n-1)$-го порядка включительно.


Это пространство изоморфно прямому произведению $L[a^k,b^k]\times R^n$,


где $L[a^k,b^k]$ --- пространство суммируемых функций


$z:[a^k,b^k]\to R^1$,


$\displaystyle \|z\|_{L[a^k,b^k]}=\int_{a^k}^{b^k} |z(t)|\,dt$.


Изоморфизм


${\cal J}=\{\Lambda_k,Y_k\}:L[a^k,b^k]\times R^n\to W^n[a^k,b^k]$


можно установить на основе представления


\begin{equation}


x(t) = \int_{a^k}^t\frac{(t-s)^{n-1}}{(n-1)!}x^{(n)}(s)\,ds+\sum_{i=0}^{n-1}


\frac{(t-a^k)^i}{i!} x^{(i)}(a^k)


\end{equation}


элемента $x\in W^n[a^k,b^k]$. Таким образом,


$$


(\Lambda_kz)(t) = \int_{a^k}^t\frac{(t-s)^{n-1}}{(n-1)!} z(s)\,ds,\quad


Y_k\beta = \sum_{i=0}^{n-1} \frac{(t-a^k)^i}{i!} \beta^i,\quad


\beta=\col \{\beta^0,\dotsc,\beta^{n-1}\};


$$


${\cal J}^{-1}=[\delta_k,r_k]$, где


$(\delta_k x)(t) = x^{(n)}(t)$,


$r_k x = \col\{x(a^k),\dotsc,x^{(n-1)}(a^k)\}$.


Для определенности будем использовать следующую норму в пространстве


$R^n$:


$|\alpha|=\sum\limits_{i=1}^n|\alpha^i|$,


$\alpha = \col\{\alpha^1,\dotsc,\alpha^n\}$. Тогда


$$


\|x\|_{W^n[a^k,b^k]}=


\|x^{(n)}\|_{L[a^k,b^k]}+\sum_{i=0}^{n-1} |x^{(i)}(a^k)|.


$$





Рассмотрим последовательность краевых задач


\begin{equation}


\lkxfk,\quad \lkxak,\quad \kodi,


\end{equation}


и ``предельную'' задачу


$$


{\cal L}_0 x=f_0,\quad l_0x=\alpha_0,


$$


где ${\cal L}_k:W^n[a^k,b^k]\to L[a^k,b^k]$


--- линейные ограниченные н\"eтеровы операторы,


$\ind {\cal L}_k=n$; $l_k:W^n[a^k,b^k]\to R^n$


--- линейные ограниченные вектор-функционалы, $k=0,1,\dotsc$





Положим $D_k=W^n[a^k,b^k]$, $B_k=L[a^k,b^k]$, $k=0,1,\dotsc$,


$c_k=\frac{b^0-a^0}{b^k-a^k}$,


$$


w_0(t) = t,\quad w_k(t)=c_k(t-a^k)+a^0,


\quad k=1,2,\dotsc


$$


Функция $w_k$ отображает отрезок $[a^k,b^k]$ на $[a^0,b^0]$;


$w^{-1}_k(t) = c^{-1}_k(t-a^0)+a^k$. Определим оператор


${\cal H}_k:B_0\to B_k$ равенством


\begin{equation}


({\cal H}_kz)(t) = z[w_k(t)].


\end{equation}


Имеем


$$


\|{\cal H}_kz\|_{B_k}=\int_{a^k}^{b^k}|z[w_k(t)]|\,dt =


c^{-1}_k\int_{a^0}^{b^0} |z(t)|\,dt.


$$


Отсюда


$\lim\limits_{k\to\infty}\|{\cal H}_kz\|_{B_k} = \|z\|_{B_0}$,


$\|{\cal H}_k\|=c^{-1}_k$,


$$


({\cal H}_k^{-1}z)(t) = z[w^{-1}_k(t)],\quad


\|{\cal H}^{-1}_k\|=c_k.


$$


Таким образом, система изоморфизмов


${\cal H} = \{{\cal H}_k\}_{k\in \N}$


является связывающей для пространств $B_0$ и $B_k$


и удовлетворяет требованию a) теоремы 1 из [1]. Сходимость


$z_k\sh z_0$ означает, что


$\displaystyle\lim_{k\to\infty}\int_{a^k}^{b^k} |z_k(t)-z_0[w_k(t)]|\,dt=0$.





Связывающую систему ${\cal P} = \{{\cal P}_k\}_{k\in \N}$


для пространств $D_0$ и $D_k$ определим, положив


\begin{equation}


({\cal P}_kx)(t) = x[w_k(t)],\quad \kodi


\end{equation}


Тогда $({\cal P}_k^{-1}x)(t)= x[w_k^{-1}(t)]$.


Имеем


$$


\frac{d^i}{dt^i}x[w_k(t)]=c^i_k


x^{(i)}[w_k(t)],\quad i=0,1,\dotsc,


$$


где


$\displaystyle


x^{(i)}[w_k(t)]=\frac{d^i}{ds^i} x(s)\Big|_{s=w_k(t)}$.


Поэтому


\begin{align*}


(\delta_k x[w_k(\cdot)])(t) &= c^n_k


x^{(n)}[w_k(t)], \\


%


\|\delta_kx[w_k(\cdot)]\|_{B_k} &= c^n_k


\int_{a^k}^{b^k} |x^{(n)}[w_k(t)]|\,dt = c^{n-1}_k


\|x^{(n)}\|_{B_0}, \\


%


r_k x[w_k(\cdot)] &= \col\{ x(a^0),


c_k\Dot{x}(a^0),\dotsc,


c^{n-1}_kx^{(n-1)}(a^0)\}, \\


%


\|{\cal P}_k x\|_{D_k} &=


c^{n-1}_k \|x^{(n)}\|_{B_0}


+ \sum_{i=0}^{n-1} c^i_k


|x^{(i)}(a^0)|.


\end{align*}


Отсюда


\begin{gather*}


\lim_{k\to\infty}\|{\cal P}_kx\|_{D_k} = \|x\|_{D_0},\\


%


\|{\cal P}_kx\|_{D_k} \leq \|x\|_{D_0} 


\max \{1, c^{n-1}_k\},


\end{gather*}


и т.\,к. равенство достигается, то


$$


\|{\cal P}_k\|=\max \{1, c^{n-1}_k\}.


$$


Точно так же


$$


\|{\cal P}_k^{-1}\| =\max \{1, c^{n-1}_k\}.


$$





Итак, система ${\cal P} = \{{\cal P}_k\}_{k\in \N}$


является связывающей для $D_0$ и $D_k$


и удовлетворяет требованию b) теоремы 1 из [1]. Сходимость


$u_k\ssp u_0$


($\lim\limits_{k\to\infty}\|u_k-u_0[w_k(\cdot)]\|_{D_k} =0$)


означает, что


\begin{equation}


\lim_{k\to\infty}\int_{a^k}^{b^k} | u_k^{(n)}(t) -


c^n_k u_0^{(n)} [w_k(t)]|\,dt = 0


\end{equation}


и


\begin{equation}


\lim_{k\to\infty} | u_k^{(i)}(a^k) -


c^i_k u_0^{(i)}(a^0)| = 0,


\quad i=0,\dotsc,n-1.


\end{equation}





Таким образом, для последовательности задач (2) справедлива


теорема 1 из [1], если $\{{\cal L}_k\}$ и $\{l_k\}$


удовлетворяют условию c) этой теоремы.


При этом уравнение


${\cal H}^{-1}_k{\cal L}_k{\cal P}_k x= f$


в задаче (7) [1] имеет вид


\begin{equation}


({\cal L}_k[x(w_k(\cdot))])(w^{-1}_k(t))=f(t).


\end{equation}





\begin{note} Система операторов $\{{\cal H}_k\}_{k\in \N}$


и $\{{\cal P}_k\}_{k\in \N}$,


определенных равенствами (3) и (4), удовлетворяет требованию (9) [1],


т.\,к. предельные соотношения (5) и (6) эквивалентны тому, что


$$


u_k^{(n)} \sh u_0^{(n)}\ \; \text{и}\ \; r_ku_k \longrightarrow r_0u_0.


$$


\end{note}





Пусть оператор ${\cal L}_k$ имеет вид


\begin{equation}


({\cal L}_k x)(t) = x^{(n)}(t) +(T_k x)(t),


\end{equation}


где $T_k:D_k\to B_k$


--- вполне непрерывный оператор. Такой вид оператор ${\cal L}_k$ имеет для


разрешенных относительно старшей производной обыкновенных дифференциальных


уравнений, уравнений с отклоняющимся аргументом, интегродифференциальных


уравнений. Оператор


$\displaystyle \frac{d^n}{dt^n}:D_k\to B_k$\;


$n$-кратного дифференцирования --- непрерывный н\"eтеров оператор,


$\displaystyle \ind\frac{d^n}{dt^n} = n$.


Поэтому оператор ${\cal L}_k:D_k\to B_k$,


определенный равенством (8), н\"eтеров,


$\ind {\cal L}_k = n$ ([2], cc.\,45, 60).





Как было отмечено выше, из сходимости $u_k \ssp u_0$


следует $u_k^{(n)}\sh u_0^{(n)}$.


Поэтому сходимость ${\cal L}_k \sph {\cal L}_0$


обеспечивается сходимостью $T_k\sph T_0$.





Так как


$$


\frac{d^n}{dt^n} ({\cal P}_kx)(t) = c^n_k x^{(n)}[w_k(t)],


$$


то уравнение (7) в этом случае имеет вид


$$


c^n_k x^{(n)}(t) +(T_k[x(w_k(\cdot))]) (w_k^{-1}(t)) = f(t).


$$





{\bf Пример.} Рассмотрим краевые задачи


\begin{gather}


x^{(n)}(t) + p_k(t)x[h_k(t)] = f_k(t),\quad t\in[a^k,b^k],\notag \\


%


x(\xi) = 0,\ \;\text{если}\ \; \xi\notin[a^k,b^k],\\


%


l_k^1x=\alpha^1,\dotsc , l_k^nx=\alpha_k^n,\quad k=0,1,\dotsc,\notag


\end{gather}


где $p_k,\,f_k\in B_k=L[a^k,b^k]$,


функция $h_k:[a^k,b^k]\to R^1$ измерима, $l_k^i$, $i=1,\dotsc,n$,


--- линейные ограниченные функционалы на пространстве


$D_k=W^n[a^k,b^k]$.





Пусть $\nu(t)$ --- определенная на отрезке $[a,b]$


вещественная функция. Для функции $x:[a,b]\to R^1$


определим операцию $S_{\nu}$ равенством


$$


(S_{\nu}x)(t) =


\begin{cases}


x[\nu(t)],& \text{если }\nu(t)\in[a,b];\\


0,& \text{если }\nu(t)\notin[a,b].


\end{cases}


$$


Используя это обозначение и полагая $l_k=[l_k^1,\dotsc,l_k^n]$,


$\alpha_k=\{\alpha_k^1,\dotsc,\alpha_k^n\}$,


запишем задачу (9) в виде


\begin{equation}


\begin{array}{c}


({\cal L}_kx)(t)\eqdef x^{(n)}(t) +


p_k(t)(S_{h_k}x)(t) = f_k(t),\\[2mm]


l_kx=\alpha_k.


\end{array}


\tag{$9'$}


\end{equation}


Полная непрерывность оператора $T_k:D_k\to B_k$,


определенного равенством


$$


(T_kx)(t) = p_k(t)(S_{h_k}x)(t),


$$


показана в ([3], c.\,17--28).





Найдем представление оператора


${\cal H}_k^{-1} T_k{\cal P}_k:D_0\to B_0$.


Пусть $x\in D_0$. Имеем


$$


(S_{h_k}{\cal P}_k x)(t) = (S_{h_k} x[w_k(\cdot)])(t) =


\begin{cases}


x[w_k(h_k(t))],& \text{если } h_k(t)\in[a^k,b^k];\\


0,& \text{если }h_k(t)\not\in[a^k,b^k].


\end{cases}


$$


Так как условие $h_k(t)\in[a^k,b^k]$ эквивалентно


$w_k(h_k(t))\in[a^0,b^0]$, то


$$


(S_{h_k}{\cal P}_k x)(t) =


\begin{cases}


x[w_k(h_k(t))],& \text{если } w_k(h_k(t))\in[a^0,b^0];\\


0,& \text{если } w_k(h_k(t))\not\in[a^0,b^0].


\end{cases}


$$


Далее,


$$


({\cal H}^{-1}_kS_{h_k}{\cal P}_k x)(t) =


\begin{cases}


x[w_k(h_k(w_k^{-1}(t)))],& \text{если } w_k(h_k(w_k^{-1}(t)))\in[a^0,b^0];\\


0,& \text{если } w_k(h_k(w_k^{-1}(t)))\notin[a^0,b^0].


\end{cases}


$$


Таким образом,


$$


({\cal H}_k^{-1} S_{h_k}{\cal P}_k\,x)(t) = (S_{g_k}x)(t),


$$


где функция $g_k:[a^0,b^0]\to R^1$ определяется равенством


$$


g_k(t) = w_k[h_k(w_k^{-1}(t))],


$$


и, следовательно,


$$


({\cal H}_k^{-1} T_k{\cal P}_kx)(t) =


(T_k[x(w_k(\cdot))])


(w_k^{-1}(t)) = p_k[w_k^{-1}(t)](S_{g_k}x)(t).


$$





Теорема 1 из [1] применительно к задаче ($9'$)


получает следующую формулировку.





{\bf Теорема 1-{\it bis}.} {\it


Пусть операторы ${\cal P}_k$, $\kodi$, и


${\cal H}_k$, $\kodi$,


определены равенствами $(4)$ и $(3)$ соответственно. Пусть, далее,


$T_k\sph T_0$ и $l_ku_k\longrightarrow l_0u_0$, если


$u_k \ssp u_0$, и задача $(9')$ при $k=0$


имеет единственное решение $x_0\in D_0$.


Для того чтобы задачи $(9')$ были однозначно разрешимы для всех достаточно


больших $k$ и для любых сходящихся последовательностей


$f_k\sh f_0$ и $\alpha_k\longrightarrow\alpha_0$ для решений


$x_k\in D_k$ этих задач имела место сходимость $x_k\ssp x_0$,


необходимо и достаточно существования такого ограниченного вектор-функционала


$l:D_0\to R^n$, что краевые задачи


\begin{equation}


\begin{array}{c}


c^n_k x^{(n)}(t) +


p_k[w_k^{-1}(t)] (S_{g_k}x)(t) = f(t),\quad t\in[a^0,b^0],\\[2mm]


lx=\alpha,\quad k=0,1,\dotsc,


\end{array}


\end{equation}


однозначно разрешимы при $k=0$ и всех достаточно больших


$k$, и при любых $f\in B_0$, $\alpha\in R^n$ для решений


$v_k\in D_0$ задач $(10)$ имеет место сходимость


$\lim\limits_{k\to\infty} \|v_k-v_0\|_{D_0} =0$.


}





\begin{note}


В приведенном примере было рассмотрено уравнение с отклоняющимся аргументом


простейшего вида с целью избежать громоздкости в преобразованиях. Та же


схема сохранится и в случае оператора


${\cal L}_k:W^n[a^k,b^k]\to L[a^k,b^k]$


более общего вида. Например, если


$$


({\cal L}_k x)(t) = x^{(n)}(t) +


\sum_{i=0}^{n-1} p_k^i(t) (S_{h_k^i} \,x^{(i)})(t),


$$


то уравнение (7) будет иметь вид


$$


c^n_k x^{(n)}(t) + \sum_{i=0}^{n-1}


p_k^i[w_k^{-1}(t)](S_{g_k^i}\;x^{(i)})(t) = f(t),


$$


где


$g_k^i(t) = w_k[h_k^i(w_k^{-1}(t))],\quad i=0,\dotsc,n-1$.


\end{note}





\begin{center}


{\bf 2. Условия сходимости


$T_k \overset{\cal P\cal H}{\longrightarrow}{} T_0$}


\end{center}


\bigskip


\addtocounter{section}{1}





Обозначим $\Omega_k = \{t\in[a^k,b^k]:h_k(t)\in [a^k,b^k]\}$;


$w_k(\Omega_k)$ --- образ множества $\Omega_k$ при отображении


$w_k$:


$w_k(\Omega_k) = \{t\in[a^0,b^0]: h_k(w_k^{-1}(t))\in[a^k,b^k]\}$;


$e_k$ --- симметрическая разность множеств $\Omega_0$ и


$w_k(\Omega_k)$:


$e_k=[\Omega_0\setminus w_k(\Omega_k)]\cup[w_k(\Omega_k)\setminus\Omega_0]$.


Таким образом, если $t\in e_k$, то либо


$h_0(t) \in[a^0,b^0]$, $h_k(w_k^{-1}(t))\notin[a^0,b^0]$,


либо $h_0(t)\not\in[a^0,b^0]$, $h_k(w_k^{-1}(t))\in[a^0,b^0]$.





{\bf Теорема.} {\it


Пусть выполнены условия





$\mathrm{a)}$ $p_k\sh p_0;$





$\mathrm{b)}$ $\lim\limits_{k\to\infty} \mes e_k =0;$





$\mathrm{c)}$ для любых $\delta >0$ и $\sigma >0$ найдется такое


$k_0$, что при $k>k_0$


\begin{equation}


\mes \{t\in\Omega_0\cup w_k(\Omega_k):


|w_k[h_k(w_k^{-1}(t))]-h_0(t)|


\geq \delta\} <\sigma.


\end{equation}


Тогда $T_k\sph T_0$.


}





\begin{note} Условие c) (см. [1]) в случае, когда все пространства


$D_k$ совпадают с $D_0$ и $h_k(t)\in[a^0,b^0]$ для


$t\in[a^0,b^0]$, $k=0,1,\dotsc$,


означает сходимость по мере последовательности


$\{h_k(t)\}$ к $h_0(t)$ на $[a^0,b^0]$.


\end{note}





\begin{pf} Из представления (1) элемента $x$


пространства $D_k$ получим оценку


$$


\max_{t\in[a^k,b^k]} |x(t)|\leq {\cal M}\|x\|_{D_k},


$$


где ${\cal M} = \max


\big\{ \frac{(b^k-a^k)^i}{i!},\ i=0,\dotsc,n-1;\ k=0,1,\dotsc\big\}.$


Если $u_k\ssp u_0$, то


$$


\sup_{k} \|u_k\|_{D_k} = {\cal M}_1 <\infty


$$


и, следовательно,


$$


\max_{t\in[a^k,b^k]} |u_k(t)| \leq {\cal M}{\cal M}_1,\quad k=0,1,\dotsc


$$





Пусть $u_k\ssp u_0$. Имеем


\begin{gather*}


(S_{h_k} u_k)(t) =


\begin{cases}


u_k[h_k(t)],& \text{если } t\in\Omega_k;\\


0,& \text{если } t\in[a^k,b^k]\setminus\Omega_k,


\end{cases}


\\


({\cal H}_k T_0 u_0)(t) = p_0[w_k(t)](S_{h_0}u_0)(w_k(t)).


\end{gather*}


Далее,


\begin{multline*}


\|T_ku_k-{\cal H}_kT_0u_0\|_{B_k} =


\int_{a^k}^{b^k}


| p_k(t) (S_{h_k}u_k)(t) -p_0[w_k(t)] (S_{h_0}u_0)(w_k(t))| \,dt\leq \\


%


\leq\int_{a^k}^{b^k}


|p_k(t)  -p_0[w_k(t)]|\, |(S_{h_k}u_k)(t)|\,dt


+ \int_{a^k}^{b^k}


|p_0[w_k(t)]|\, |(S_{h_k}u_k)(t) -(S_{h_0}u_0)(w_k(t))|\,dt.


\end{multline*}


Первое слагаемое в правой части неравенства оценивается величиной


$$


{\cal M}{\cal M}_1 \int_{a^k}^{b^k}


|p_k(t) -p_0[w_k(t)] |\,dt,


$$


следовательно, стремится к нулю при $k\to\infty$


в силу условия a).





Второе слагаемое равно


$$


c^{-1}_k\int_{a^0}^{b^0}


|p_0(s)| |(S_{h_k}u_k)(w_k^{-1}(s)) - (S_{h_0}u_0)(s)|\,ds \eqdef


c^{-1}_k{\cal I}_{[a^0,b^0]}.


$$


Интеграл ${\cal I}_{[a^0,b^0]}$ разобьем на два слагаемых


$$


{\cal I}_{[a^0,b^0]} =


{\cal I}_{\Omega_0\cap w_k(\Omega_k)} + {\cal I}_{e_k}.


$$


Если $t\in e_k$, то либо


$(S_{h_k}u_k)(w_k^{-1}(t)) = u_k[h_k(w_k^{-1}(t))]$,


$(S_{h_0}u_0)(t) =0$,


либо $(S_{h_k}u_k)(w_k^{-1}(t)) = 0$,


$(S_{h_0}u_0)(t) =u_0[h_0(t)]$.


Поэтому имеем оценку


$$


{\cal I}_{e_k} \leq {\cal M}{\cal M}_1 \int_{e_k} |p_0(s)|\,ds,


$$


откуда


$\lim\limits_{k\to\infty} {\cal I}_{e_k} = 0$


в силу условия b).


Далее,


\begin{align*}


{\cal I}_{\Omega_0\cap w_k(\Omega_k)}&\leq


\int_{\Omega_0\cap w_k(\Omega_k)}


|p_0(s)|\,\big |u_k[h_k(w_k^{-1}(s))] -


u_0[w_k(h_k(w_k^{-1}(s)))]\big|\,ds + \\


%


&\qquad+\int_{\Omega_0\cap w_k(\Omega_k)}


|p_0(s)|\,\big |u_0[w_k(h_k(w_k^{-1}(s)))] -


u_0[h_0(s)]\big|\,ds \eqdef {\cal I}_{1k} + {\cal I}_{2k},\\


%


{\cal I}_{1k} &\leq {\cal M} \|u_k-u_0(w_k)\|_{D_k}


\int_{a^0}^{b^0} |p_0(s)|\,ds,


\end{align*}


откуда $\lim\limits_{k\to\infty} {\cal I}_{1k} =0$.





Пусть задано $\varepsilon>0$. Найдем $\delta >0$ и


$\sigma >0$ так, чтобы


$$


|u_0(t_1)-u_0(t_2)|<\varepsilon,\ \; \text{если}\ \;|t_1-t_2|<\delta,


\quad t_1,t_2\in[a^0,b^0],


$$


и


$\displaystyle \int_{e}|p_0(s)|\,ds <\varepsilon$,


если $\mes e<\sigma$. Выберем $k_0$ так, чтобы при $k>k_0$


выполнялось неравенство (11). Обозначим


\begin{align*}


E_{1k} &= \{s\in\Omega_0\cap w_k(\Omega_k):


|w_k[h_k(w_k^{-1}(s))]-h_0(s)|


\geq \delta\},\\


%


E_{2k} &= (\Omega_0\cap w_k(\Omega_k)) \setminus E_{1k}.


\end{align*}


Тогда


\begin{multline*}


{\cal I}_{2k} =


\int_{E_{1k}} |p_0(s)|\,\big |u_0[w_k(h_k(w_k^{-1}(s)))] -


u_0[h_0(s)]\big|\,ds +\\


%


+\int_{E_{2k}} |p_0(s)|\,\big |u_0[w_k(h_k(w_k^{-1}(s)))] -


u_0[h_0(s)]\big|\,ds \leq \\


%


\leq 2{\cal M}{\cal M}_1 \int_{E_{1k}} |p_0(s)|\,ds +


\varepsilon \int_{a^0}^{b^0}


|p_0(s)|\,ds < \varepsilon (2{\cal M}{\cal M}_1 + \|p_0\|_{B_0})


\end{multline*}


для $k>k_0$. Итак,


$\lim\limits_{k\to\infty} {\cal I}_{2k} =0$.


\end{pf}





Теорема для случая, когда все пространства $D_k$ совпадают с


$D_0$, была сформулирована в [4].





\section{Уравнение $n$-го порядка с импульсными воздействиями}





Обозначим через $\wsn$ пространство функций $x:[a,b]\to R^1$,


представимых в виде


$$


x(t) = \int_a^t \frac{(t-s)^{n-1}}{(n-1)!} z(s)\,ds + \sum_{i=0}^{n-1}


\frac{(t-a)^i}{i!} \beta_i +


\sum_{j=1}^m \chi_{[t_j,b]}(t) \frac{(t-t_j)^{n-1}}{(n-1)!} \gamma_j,


$$


где $z\in L[a,b]$;


$\beta_0,\dotsc,\beta_{n-1}$, $\gamma_1,\dotsc,\gamma_m$


--- постоянные; $\chi_{[t_j,b]}$


--- характеристическая функция отрезка $[t_j,b]$,


$a<t_1<\dotsb<t_m<b$.


Таким образом, элементами пространства $\wsn$


являются функции, имеющие на $[a,b]$


абсолютно непрерывные производные до


$(n-2)$-го порядка включительно (если $n\geq 2$) и производную


$(n-1)$-го порядка, абсолютно непрерывную на каждом из промежутков


$[a,t_1), [t_1,t_2), \dotsc, [t_m,b]$.





Пространство $\wsn$ изоморфно прямому произведению


$L[a,b]\times R^{n+m}$. Изоморфизм $\{\Lambda,Y\}$


определяется равенствами


\begin{gather*}


(\Lambda z)(t) = \int_a^t \frac{(t-s)^{n-1}}{(n-1)!} z(s)\,ds,\\


%


(Y\beta)(t) = \sum_{i=0}^{n-1} \frac{(t-a)^i}{i!} \beta_i + \sum_{j=1}^m


\chi_{[t_j,b]}(t) \frac{(t-t_j)^{n-1}}{(n-1)!} \gamma_j,\\


%


\beta= \col\{\beta_0,\dotsc, \beta_{n-1},\; \gamma_1,\dotsc,\gamma_m\};


\end{gather*}


$\{\Lambda,Y\}^{-1} = [\delta,r]$, где


$$


(\delta x)(t) = x^{(n)} (t),\quad


rx = \col \{x(a),\dotsc,x^{(n-1)}(a),


\Delta x^{(n-1)}(t_1),\dotsc,\Delta x^{(n-1)}(t_m)\},


$$


$\Delta x^{(n-1)}(t_i) = x^{(n-1)}(t_i) - x^{(n-1)}(t_i-0)$;


$$


\|x\|_{\wsn}=\|x^{(n)}\|_{L[a,b]} + |rx|.


$$





Пусть


$\lim\limits_{k\to\infty} a^k = a^0$,


$\lim\limits_{k\to\infty} b^k=b^0$,


$\lim\limits_{k\to\infty} t_i^k = t_i^0$, $i=1,\dotsc,m$.


Обозначим


$$


B_k=L[a^k,b^k],\quad D_k=WS^n[a^k,t_1^k,\dotsc,t_m^k,b^k],\quad k=0,1,\dotsc


$$


Пространство $D_k$ изоморфно $B_k\times R^{n+m}$. Изоморфизмы


$$


\{\Lambda_k,Y_k\}:B_k\times R^{n+m} \to D_k


\ \; \text{и}\ \;


[\delta_k,r_k]= \{\Lambda_k,Y_k\}^{-1}:D_k\to B_k\times R^{n+m}


$$


определяются равенствами


\begin{gather*}


(\Lambda_k z)(t) = \int_{a^k}^t \frac{(t-s)^{n-1}}{(n-1)!} z(s)\,ds,\\


%


(Y_k \beta)(t)=\sum_{i=0}^{n-1}\frac{(t-a^k)^i}{i!} \beta_i + \sum_{j=1}^m


\chi_{[t_j^k, b^k]}(t) \frac{(t-t_j^k)^{n-1}}{(n-1)!} \gamma_j, \\


%


(\delta_k x)(t) = x^{(n)} (t), \\


%


r_k x = \col \{x(a^k),\dotsc, x^{(n-1)}(a^k), \Delta x^{(n-1)}(t_1^k),


\dotsc,\Delta x^{(n-1)}(t_m^k)\}.


\end{gather*}


Конечномерный оператор $Y_k:R^{n+m}\to D_k$ определяется вектором


$$


Y_k(t)=\bigg\{1,\ \frac{t-a^k}{1!}, \dotsc,\frac{(t-a^k)^{n-1}}{(n-1)!},


\ \chi_{[t_1^k, b^k]}(t)\frac{(t-t_1^k)^{n-1}}{(n-1)!},\dotsc,


\chi_{[t_m^k, b^k]}(t)\frac{(t-t_m^k)^{n-1}}{(n-1)!}\bigg\}.


$$


Таким образом, элемент $x\in D_k$, $k=0,1,\dotsc$, имеет представление


\begin{multline}


x(t) = (\Lambda_k x^{(n)})(t) + Y_k (t)r_kx =


\int_{a^k}^t \frac{(t-s)^{n-1}}{(n-1)!} x^{(n)}(s)\,ds +\\


%


+ \sum_{i=0}^{n-1}\frac{(t-a^k)^i}{i!} x^{(i)}(a^k)


+ \sum_{j=1}^m \chi_{[t_j^k, b^k]}(t)


\frac{(t-t_j^k)^{n-1}}{(n-1)!} \Delta x^{(n-1)}(t_j^k);


\end{multline}


$\|x\|_{D_k}=\|\delta_kx\|_{B_k} + |r_kx|$.





Положим $w_0(t) = t$,


\begin{equation}


w_k(t) = \sum_{i=0}^m


\bigg[\frac{t^0_{i+1}-t^0_i}{t^k_{i+1}-t^k_i} (t-t^k_i) + t^0_i\bigg]


\chi_{[t^k_i,t^k_{i+1}]}(t),\quad k=1,2,\dotsc


\end{equation}


(здесь $t_0^k=a^k$, $t^k_{m+1}=b^k$). Функция $w_k$


отображает отрезок $[a^k,b^k]$ на $[a^0,b^0]$ так, что


$w_k(t^k_i)=t^0_i$, $i=0,\dotsc,m+1$.





Функции $w_k(t)$,


определяемые равенством (13), были использованы А.В.\,Анохиным при изучении


последовательности краевых задач для систем функционально-дифференциальных


уравнений первого порядка с импульсными воздействиями. Соответствующий


результат опубликован без доказательства в ([3], теорема 6.3.2).





Определим оператор ${\cal H}_k:B_0\to B_k$, $k=0,1,\dotsc $,


равенством


\begin{equation}


({\cal H}_k z)(t) = z[w_k(t)].


\end{equation}


Тогда $({\cal H}_k^{-1} z)(t) = z[w_k^{-1}(t)]$,


$$


w_k^{-1}(t)=\sum_{i=0}^m\bigg[\frac{t^k_{i+1}-t^k_i}{t^0_{i+1}-t^0_i}


(t-t^0_i)+t^k_i\bigg] \chi_{[t^0_i,t^0_{i+1}]}(t).


$$


Для $z\in D_0$ имеем


$$


\|{\cal H}_k z\|_{B_k} = \int_{a^k}^{b^k} |z[w_k(t)]|\,dt


= \sum_{i=0}^m\int_{t^k_i}^{t^k_{i+1}}


|z[w_k(t)]|\,dt =


\sum_{i=0}^m\frac{t^k_{i+1}-t^k_i}{t^0_{i+1}-t^0_i}


\int_{t^0_i}^{t^0_{i+1}}|z(\tau)|\,d\tau.


$$


Отсюда


$\lim\limits_{k\to\infty} \|{\cal H}_kz\|_{B_k} = \|z\|_{B_0}$,


$\|{\cal H}_k z\|_{B_k} \leq


\displaystyle\max_i\dfrac{t^k_{i+1}-t^k_i}{t^0_{i+1}-t^0_i}\,


\|z\|_{B_0}$.


Так как равенство достигается, то


$$


\|{\cal H}_k \| = \max_i\frac{t^k_{i+1}-t^k_i}{t^0_{i+1}-t^0_i}.


$$


Аналогично


$$


\|{\cal H}^{-1}_k\| = \max_i\frac{t^0_{i+1}-t^0_i}{t^k_{i+1}-t^k_i}.


$$


Итак, система изоморфизмов


${\cal H} = \{{\cal H}_k\}_{k\in \N}$


является связывающей для $B_0$ и $B_k$, причем


$\sup\limits_k\|{\cal H}_k^{-1}\| < \infty$.


Сходимость $z_k\sh z_0$ означает, что


$$


\lim_{k\to\infty} \int_{a^k}^{b^k}


|z_k(t) - z_0[w_k(t)]|\,dt =0.


$$





Связывающую систему ${\cal P} = \{{\cal P}_k\}_{k\in \N}$


изоморфизмов для пространств $D_0$ и $D_k$


построим так, как это было сделано в \S\,4 из [1].


А именно, положим


\begin{equation}


{\cal P}_k=\Lambda_k{\cal H}_k\delta_0 + Y_k r_0.


\end{equation}


Таким образом,


\begin{multline}


({\cal P}_kx)(t) =


\int_{a^k}^t \frac{(t-s)^{n-1}}{(n-1)!} x^{(n)}[w_k(s)]\,ds +


\sum_{i=0}^{n-1} \frac{(t-a^k)^i}{i!} x^{(i)}(a^0)+\\


%


+\sum_{j=1}^m \frac{(t-t_j^k)^{n-1}}{(n-1)!} \chi_{[t_j^k,b^k]}(t)


\Delta x^{(n-1)}(t_j^0),\quad x\in D_0.


\end{multline}


Для $u_k\in D_k$, $k=0,1,\dotsc$, имеем


$$


\|u_k - {\cal P}_k u_0\|_{D_k} =


\int_{a^k}^{b^k} |u_k^{(n)}(s) - u_0^{(n)}[w_k(s)]|\,ds +


\sum_{i=0}^{n-1} |u_k^{(i)}(a^k)-u^{(i)}_0(a^0)| +


\sum_{j=1}^m |\Delta u_k^{(n-1)}(t_j^k) - \Delta u_0^{(n-1)}(t_j^0)|.


$$


Поэтому сходимость $u_k\ssp u_0$ означает, что


$u_k^{(n)}\sh u_0^{(n)}$ и $r_ku_k\to r_0u_0$.





Пусть ${\cal L}_k:D_k\to B_k$


--- линейный ограниченный н\"eтеров оператор,


$\ind {\cal L}_k = n+m$. Применяя ${\cal L}_k$


к обеим частям равенства (16), получим


$({\cal L}_k{\cal P}_kx)(t) = (Q_k x^{(n)}[w_k(\cdot)])(t) + A_k(t) r_0x$,


где $Q_k={\cal L}_k\Lambda_k$ --- ``главная часть'' оператора


${\cal L}_k$ ([3], c.\,103),


$A_k(t) = \{ A_0^k(t),\dotsc,A^k_{n-1}(t), B^k_1(t),\dotsc,B_m^k(t)\} =


({\cal L}_k Y_k(\cdot ))(t)$.





Связывающие системы изоморфизмов пространств


$B_0$, $B_k$ и $D_0$, $D_k$, определенные равенствами (14) и (15)


соответственно, удовлетворяют требованиям


теоремы 1. Поэтому для краевых задач


$$


{\cal L}_k x= f_k,\quad l_k x=\alpha_k\in R^{n+m},\quad k=0,1,\dotsc,


$$


в предположении, что ${\cal L}_k \sph {\cal L}_0$


и $l_ku_k\to l_0u_0$, если $u_k\ssp u_0$,


справедлива теорема 1. При этом уравнение


${\cal H}^{-1} {\cal L}_k{\cal P}_kx=f$


в задаче (7) из [1] имеет вид


$$


(Q_kx^{(n)}[w_k(\cdot)])(w_k^{-1}(t)) +


\sum_{i=0}^{n-1} A_i^k(w_k^{-1}(t)) x^{(i)}(a^0)+


\sum_{j=1}^m B_j^k(w_k^{-1}(t))\Delta x^{(n-1)}(t_j^0) = f(t).


$$





Применяя оператор ${\cal L}_k$


к обеим частям равенства (12), получим


$({\cal L}_k x)(t) = (Q_kx^{(n)})(t) + A_k(t) r_kx$.


Из равенства


${\cal L}_k u_k - {\cal H}_k{\cal L}_0 u_0 =


Q_k u_k^{(n)}-{\cal H}_k Q_0 u_0^{(n)}+A_k r_k u_k-{\cal H}_k A_0 r_0 u_0$


и оценки


$\|{\cal L}_k u_k-{\cal H}_k{\cal L}_0 u_0\|_{B_k} \leq


\|Q_k u^{(n)}_k-{\cal H}_k Q_0u^{(n)}_0\|_{B_k} +


\|\,|A_k-{\cal H}_k A_0|\,\|_{B_k}


|r_ku_k| +\|\,|{\cal H}_k A_0|\,\|_{B_k} |r_ku_k-r_0u_0|$


следует, что сходимость ${\cal L}_k\sph{\cal L}_0$


эквивалентна тому, что $Q_k\shh Q_0$ и


$A^k_i\sh A^0_i$, $i=0,\dotsc,n-1$,


$B_j^k\sh B_j^0$, $j=1,\dotsc,m$.





\begin{note}


В случае уравнения первого порядка $(n=1)$


можно связывающую систему для $D_0$ и $D_k$


определить равенством


$({\cal P}_k x)(t)=x[w_k(t)]$, $\kodi$


Тогда уравнение


${\cal H}_k^{-1} {\cal L}_k{\cal P}_k x= f$


будет иметь такой же вид, как (7),


причем $w_k$ определено равенством (13).


\end{note}





\begin{thebibliography}{99}





\bibitem{1}


Анохин А.В., Рахматуллина Л.Ф.


{\em Непрерывная зависимость от параметров


решения линейной краевой задачи.} I // Изв. вузов.


Математика. -- 1996. -- \No\,11. -- С.\,29--38.


\bibitem{2}


Крейн С.Г. {\em Линейные уравнения в банаховом 


пространстве.} -- М.: Наука, 1971. -- 104 с.


\bibitem{3}


Азбелев Н.В., Максимов В.П., Рахматуллина Л.Ф. {\em Введение в теорию


функционально-дифференциальных уравнений}. -- М.: Наука, 1991. -- 280~с.


\bibitem{4}


Рахматуллина Л.Ф. {\em


О последовательности краевых задач для линейных уравнений с отклоняющимся


аргументом} // Краев. задачи. -- Пермь, 1982. -- С.\,18--22.





\end{thebibliography}





\vskip 0.5cm





\post{\it Пермский государственный}{\it Поступила}


\post{\it технический университет}{30.04.1999}





\label{end}


\end{document}


